\documentclass[12pt, a4paper]{article}
\usepackage{config}

\begin{document}

\begin{titlepage}
    \centering
    \begin{flushleft}
        \includegraphics[scale=0.35]{images/INF-LOGO.png}
    \end{flushleft}
    \vspace*{6cm}

    {\fontsize{24}{28.8}\selectfont\bfseries Informe de Práctica Inicial}\\[0.7cm]
    {\fontsize{16}{19.2}\selectfont Universidad Católica de Temuco}\\[0.3cm]
    {\fontsize{14}{16.8}\selectfont Nombre de la Organización o Empresa}

    \vfill

    \begin{tabular}{@{}ll@{}}
        \textbf{Autor:} & Nombre del estudiante en práctica\\[0.3cm]
        \textbf{Profesor encargado:} & Nombre del profesor encargado\\[0.3cm]
        \textbf{Supervisor de práctica:} & Nombre del supervisor en la empresa\\[0.3cm]
        \textbf{Carrera:} & Ingeniería Civil en Informática\\[0.3cm]
        \textbf{Fecha de inicio:} & DD de Mes de AAAA\\[0.3cm]
        \textbf{Fecha de término:} & DD de Mes de AAAA\\[0.3cm]
        \textbf{Fecha de informe:} & DD de Mes de AAAA\\
    \end{tabular}

    \vspace{2cm}
\end{titlepage}

\tableofcontents
\newpage

\section{Resumen}
Este informe detalla las características de la práctica inicial/profesional universitaria realizada por el autor, incluyendo la empresa donde se llevó a cabo, las funciones desempeñadas y las actividades desarrolladas durante el período de práctica.

La integral de Gauss es $\int_{-\infty}^{\infty} e^{-x^2} dx = \sqrt{\pi}$.

\noindent\textbf{Palabras clave:} Palabra clave 1, Palabra clave 2, Palabra clave 3, Palabra clave 4.

\section{Introducción}
Redactar el contexto general de la práctica, los objetivos generales y la estructura del informe.

\begin{figure}[h]
    \centering
    \includegraphics[width=0.8\textwidth]{images/chart.jpg}
    \caption{Gráfico de rendimiento del sistema desarrollado.}
    \label{fig:chart}
\end{figure}

\section{Objetivos y expectativas}
En esta sección deberás redactar un resumen de los objetivos y expectativas esperados al finalizar el proceso de práctica. Ejemplo como contribuirá este proceso en tu desarrollo social, humano y profesional. (Conocimientos técnicos, habilidades sociales, etc).

La ecuación de Schrödinger es:
\begin{equation}
i\hbar \frac{\partial}{\partial t} \Psi = \hat{H} \Psi
\end{equation}

\section{Descripción de la empresa}
Describir la organización, su ámbito de trabajo y relevancia en el contexto donde opera.

\subsection{Misión}
Describir el propósito fundamental y la razón de ser de la organización.

\subsection{Visión}
Describir la proyección a futuro y los objetivos estratégicos de la organización.

\subsection{Funciones}
Con respecto a las funciones que se desarrollan en la empresa/organización, se establece que se realizan las siguientes funciones generales:

\begin{itemize}
    \item Desarrollo de software
    \begin{itemize}
        \item Frontend con React
        \item Backend con Node.js
    \end{itemize}
    \item Gestión de bases de datos
    \item Soporte técnico
\end{itemize}

\section{Organigrama de la empresa}
Describir la estructura de la empresa/departamento/dirección donde se realiza la práctica. Incluir un diagrama visual del organigrama si es posible.

\subsection{Departamento o área específica}
Describir el departamento o área donde se realizó la práctica, sus funciones y cómo se relaciona con el resto de la organización.

\begin{table}[h]
    \centering
    \caption{Estadísticas de uso del sistema.}
    \label{tab:stats}
    \begin{tabular}{|l|r|r|r|}
        \hline
        \textbf{Métrica} & \textbf{Valor} & \textbf{Unidad} & \textbf{Notas} \\
        \hline
        Usuarios activos & 1,250 & personas & Promedio mensual \\
        \hline
        Tiempo de respuesta & 150 & ms & Percentil 95 \\
        \hline
        Tasa de error & 0.05 & \% & Máximo aceptable \\
        \hline
        \multicolumn{3}{|c|}{Total de sesiones} & 15,000 \\
        \hline
    \end{tabular}
\end{table}

\section{Actividades encomendadas}
Describir detalladamente las principales actividades realizadas durante la práctica. Explicar el problema a resolver, la solución implementada y los resultados obtenidos. Incluir capturas de pantalla, diagramas o ejemplos cuando sea relevante.

\subsection{Actividad 1: Desarrollo de API}
Desarrollo de una API RESTful para gestión de usuarios.

\subsection{Actividad 2: Implementación de autenticación}
Implementación de JWT para autenticación segura.

\section{Tecnologías aplicadas}
Un stack tecnológico es el conjunto de tecnologías que se utilizan para construir y ejecutar aplicaciones. Está compuesto por diferentes herramientas que interactúan entre sí para hacer que una aplicación funcione de manera eficiente. Puede declarar secciones categorizando por ejemplo: Lenguaje de programación, Base de datos, Frameworks y librerías, entre otros.

\begin{enumerate}
    \item Lenguaje: Rust para backend, TypeScript para frontend.
    \item Base de datos: MongoDB para datos no estructurados.
    \item Librerías: Serde para serialización, Tokio para async.
\end{enumerate}

\section{Experiencia en el proceso de práctica}
Reflexionar sobre los aprendizajes técnicos adquiridos, el desarrollo de habilidades blandas (comunicación, trabajo en equipo, presentación de avances), los desafíos enfrentados y los aspectos positivos de la experiencia.

\section{Reflexiones}
Realizar una reflexión crítica sobre la experiencia completa de práctica. Evaluar el impacto del trabajo realizado, plantear perspectivas futuras para el proyecto o la organización, y agradecer a las personas que apoyaron el proceso.

\newpage
\printbibliography[title={Referencias}]

\newpage
\section*{Anexos}
\addcontentsline{toc}{section}{Anexos}

\subsection*{Anexo A: Bitácoras de práctica}
Incluir las bitácoras del proceso de práctica.

\subsubsection*{Semana 1}

\begin{tabularx}{\linewidth}{|>{\raggedright\arraybackslash}m{3cm}|X|}
\hline
\textbf{Fecha} & \textbf{Actividades} \\
\hline
\rowcolor{gray!20}
07/01/2025 & Inicio del proyecto \\
\cline{1-2}
\multicolumn{2}{|X|}{Configuración inicial del entorno y revisión de requisitos.} \\
\hline
\end{tabularx}

\subsection*{Anexo B: Material complementario}
Diagramas y documentación adicional.

\subsection*{Anexo C: Repositorio}
Enlace al repositorio: https://github.com/example/repo

\end{document}